\documentclass[a4paper,11pt]{article}

% Engine-specific font set-up (the style files themselves contain none)
\usepackage{iftex}
\ifPDFTeX
  \usepackage[T1]{fontenc}
  \usepackage{lmodern}
\else
  \usepackage{fontspec}
\fi

\usepackage[british]{babel}
\usepackage{csquotes}
\usepackage[style=leeds-harvard]{biblatex}
\addbibresource{leeds-harvard-test.bib}
\usepackage[hidelinks]{hyperref}

\usepackage[margin=2.5cm]{geometry}
\usepackage{array}
\usepackage{longtable}
\usepackage{booktabs}
\setlength{\parskip}{0.5em}
\setlength{\parindent}{0pt}

\title{Leeds-Harvard Referencing Style for BibLaTeX}
\author{Charles Lam\\Language Centre, University of Leeds}
\date{\today}

\begin{document}
\maketitle

\begin{abstract}
This document demonstrates the \texttt{leeds-harvard} BibLaTeX style. It details the exact \LaTeX{} commands, rendered in-text citations, reference list formatting, and specific rules in the Leeds-Harvard style.
\end{abstract}

\section{In-Text Citations: Syntax \& Output}

Leeds-Harvard is an author--date referencing system. Table~\ref{tab:citations} details the standard citation commands alongside their rendered output:

\begingroup
\small
\begin{longtable}{@{} >{\raggedright}p{3.8cm} >{\raggedright}p{6.4cm} >{\raggedright\arraybackslash}p{4.8cm} @{}}
\caption{In-text citation syntax and rendered Leeds-Harvard output.} \label{tab:citations} \\
\toprule
\textbf{Citation Scenario} & \textbf{\LaTeX{} Command} & \textbf{Rendered Output} \\
\midrule
\endfirsthead
\toprule
\textbf{Citation Scenario} & \textbf{\LaTeX{} Command} & \textbf{Rendered Output} \\
\midrule
\endhead
\bottomrule
\endfoot
Pparenthetical & \texttt{\textbackslash parencite\{shannon1948\}} & \parencite{shannon1948} \\
Narrative & \texttt{\textbackslash textcite\{shannon1948\}} & \textcite{shannon1948} \\
Specific page & \texttt{\textbackslash parencite[24]\{shannon1948\}} & \parencite[24]{shannon1948} \\
Page range & \texttt{\textbackslash parencite[20--29]\{shannon1948\}} & \parencite[20--29]{shannon1948} \\
Roman numerals & \texttt{\textbackslash parencite[iv]\{dirac1958\}} & \parencite[iv]{dirac1958} \\
Organization author (first cite) & \texttt{\textbackslash parencite\{nasa2011\}} & \parencite{nasa2011} \\
Organization author (subsequent) & \texttt{\textbackslash parencite\{nasa2011\}} & \parencite{nasa2011} \\
Same author, multiple years & \texttt{\textbackslash parencite\{knuth1968,}\allowbreak\texttt{knuth1984texbook,}\allowbreak\texttt{knuth1997\}} & \parencite{knuth1968,knuth1984texbook,knuth1997} \\
Multiple works, different authors & \texttt{\textbackslash parencite\{macnaughton1999,}\allowbreak\texttt{ahmed2012,knuth1997\}} & \parencite{macnaughton1999,ahmed2012,knuth1997} \\
Narrative list of authors & \texttt{\textbackslash textcite\{adams1906,}\allowbreak\texttt{strogatz1994,han2026\}} & \textcite{adams1906,strogatz1994,han2026} \\
\bottomrule
\end{longtable}
\endgroup

\section{What Users Should Expect from Leeds-Harvard Referencing}

\begin{enumerate}
\item Leeds-Harvard use the English word ``and'' instead of an ampersand (``\&''), e.g. \parencite{gencturk2015}) and in the reference list (e.g. \emph{Tanenbaum, A.S. and Wetherall, D.J.}).

\item How does Leeds-Harvard handle multiple publications by the same author?
\textbf{What to expect:} Leeds-Harvard compresses author names in parenthetical citations. If you cite multiple works by Donald Knuth, the author's surname is printed once followed by the publication years separated by semicolons:
\begin{itemize}
  \item \texttt{\textbackslash parencite\{knuth1968,knuth1984texbook,knuth1997\}} \\
    yields \parencite{knuth1968,knuth1984texbook,knuth1997}.
\end{itemize}

\item In what order should multiple citations appear inside the same parentheses?}
\textbf{What to expect:} Leeds-Harvard specifies that multiple sources in one citation must be arranged \textbf{chronologically} (oldest first), separated by semicolons:
\begin{itemize}
  \item \parencite{macnaughton1999,ahmed2012,knuth1997}.
  \item \textbf{Important:} Because \texttt{sortcites=false} is set to honour the author's intentional citation order, list your BibTeX keys in chronological order in your \texttt{\textbackslash parencite} call.
\end{itemize}

\item When is ``et al.'' used in citations versus the reference list?\\
\textbf{What to expect:}
\begin{itemize}
  \item \textbf{In-Text Citations:}
    \begin{itemize}
      \item 1 or 2 authors: Both surnames should appear (\parencite{shannon1948}, \parencite{gencturk2015}).
      \item 3 or more authors: Only the first author followed by ``et al.'' (\parencite{macnaughton1999}).
      \item Note that ``et al.'' is supposed to be written in upright roman font (i.e. \textbf{not italicised}) in Leeds-Harvard. 
    \end{itemize}
  \item \textbf{Reference List:} \emph{All} authors must be explicitly named regardless of how many there are. The abbreviation ``et al.'' is \textbf{not} permitted in the bibliography itself.
\end{itemize}

\item How are page numbers, page ranges, roman numerals, and unpaginated works formatted?\\
\begin{itemize}
  \item \textbf{Single page:} Lowercase \texttt{p.} without a space: \parencite[24]{shannon1948}.
  \item \textbf{Page ranges:} \texttt{pp.} with a hyphen and no space: \parencite{shannon1948} (in bibliography: \emph{pp.379-423}).
  \item \textbf{Roman numerals:} Do \emph{not} add \texttt{p.} to roman numerals: \texttt{\textbackslash parencite[iv]\{dirac1958\}} $\to$ \parencite[iv]{dirac1958}.
  \item \textbf{No pagination:} If a work is unpaginated (e.g. an online article), specify \texttt{no pagination} in the postnote: \texttt{\textbackslash parencite[no pagination]\{dirac1958\}} $\to$ \parencite[no pagination]{dirac1958}.
\end{itemize}

\item Why do corporate or organization authors require double braces in BibTeX?\\
If a corporate name is entered with single braces:
\begin{verbatim}
author = {National Aeronautics and Space Administration}
\end{verbatim}
BibTeX parses ``Administration'' as the surname and ``National Aeronautics and Space'' as first names!
\begin{itemize}
  \item \textbf{Double braces in \texttt{.bib}:} Corporate names must be protected with double curly braces:
  \begin{verbatim}
  author      = {{National Aeronautics and Space Administration}},
  shortauthor = {NASA},
  \end{verbatim}
  \item \textbf{Automatic acronym expansion:} The first citation prints the full corporate name with the acronym in square brackets: \parencite{nasa2011}. Subsequent citations automatically use the acronym: \parencite{nasa2011}.
\end{itemize}

\item What happens when a book lacks a place of publication or publisher?
\begin{itemize}
  \item In standard Leeds-Harvard guidelines, works with an omitted place of publication historically printed \texttt{[no place]}.
  \item In modern academic practice, place of publication is often omitted for major publishers (e.g. \parencite{hawking1988}). By default, \texttt{leeds-harvard} sets \texttt{no-place=false} to cleanly omit \texttt{[no place]}.
  \item To enforce strict Leeds-Harvard behaviour (where every missing publication place prints \texttt{[no place]: Publisher}), load the package with:
  \begin{verbatim}
  \usepackage[style=leeds-harvard, no-place=true]{biblatex}
  \end{verbatim}
  \item If a book has no publisher specified, it correctly prints \texttt{[no publisher]} (e.g. Gibbs 1881).
\end{itemize}

\item Why should titles in the \texttt{.bib} file be in sentence case?\\
Leeds-Harvard requires article and book titles to appear in \textbf{sentence case} (only the first word of the title, subtitle, and proper nouns capitalized). BibLaTeX preserves title capitalization as entered in your \texttt{.bib} file. Always protect proper nouns, acronyms, and formulas with curly braces:
\begin{verbatim}
  title = {The {Byzantine} generals problem},
  title = {Evaluation of reinforced concrete and reinforced
           engineered cementitious composite ({ECC}) members},
\end{verbatim}

\item Why is there no full stop at the end of reference entries ending in a URL or DOI?\\
To avoid broken hyperlinks when readers copy text from a PDF, Leeds-Harvard strictly prohibits placing a full stop immediately after a URL or DOI. Trailing punctuation is automatically stripped by the package.





\end{enumerate}








\section{Physics-Specific Citations and Journals}

In Physics and Astronomy, referencing conventions often rely on automated bibliographic exports (e.g. SAO/NASA ADS, arXiv, APS, AIP, IOP). Leeds-Harvard accommodates these workflows while maintaining university standards:

\subsection{Full Journal Titles vs Standard Physics Macros}
Physics authors routinely use standard AAS/APS macros in BibTeX exports (such as \texttt{\textbackslash apjl}, \texttt{\textbackslash mnras}, \texttt{\textbackslash aap}, \texttt{\textbackslash prl}, \texttt{\textbackslash pra}). However, \textbf{Leeds-Harvard strictly forbids journal abbreviations}.

The \texttt{leeds-harvard} style includes an internal source map that automatically expands these macros into their full official titles during compilation:
\begin{itemize}
  \item \texttt{\textbackslash apjl} $\to$ \emph{The Astrophysical Journal Letters}
  \item \texttt{\textbackslash mnras} $\to$ \emph{Monthly Notices of the Royal Astronomical Society}
  \item \texttt{\textbackslash aap} $\to$ \emph{Astronomy and Astrophysics}
  \item \texttt{\textbackslash prl} $\to$ \emph{Physical Review Letters}
  \item \texttt{\textbackslash pra}, \texttt{\textbackslash prb}, \texttt{\textbackslash prc}, \texttt{\textbackslash prd}, \texttt{\textbackslash pre} $\to$ \emph{Physical Review A}, \emph{B}, \emph{C}, \emph{D}, \emph{E}
\end{itemize}
You do \textbf{not} need to load \texttt{aas\_macros.sty} or manually rewrite your BibTeX database.

\subsection{Electronic Article Identifiers (\texttt{eid})}
Modern physics journals (APS *Physical Review*, AIP, IOP, *ApJ Letters*) use article identifiers instead of sequential pagination:
\begin{itemize}
  \item \parencite{kurikawa2020}: \emph{Physical Review Research}, \textbf{2}(2), article no: 023307 [no pagination].
  \item \parencite{2026ApJ..1008L..60C}: \emph{The Astrophysical Journal Letters}, \textbf{1008}(2), article no: L60 [no pagination].
\end{itemize}
The package detects \texttt{eid} (or \texttt{eid = pages} from ADS) and formats it as \emph{article no: \dots{} [no pagination]}.

\subsection{arXiv Preprints}
Preprints from arXiv are fully supported with official Leeds-Harvard preprint formatting:
\begin{itemize}
  \item \textcite{2015arXiv150102297G} formats as: \\
    \emph{Grachev, A.A., Leo, L.S., Di Sabatino, S., Fernando, H.J.S., Pardyjak, E.R. and Fairall, C.W. 2015. [Pre-print]. Structure of turbulence in katabatic flows below and above the wind-speed maximum. arXiv. [Online]. Available from:} \url{https://doi.org/10.48550/arXiv.1501.02297}
\end{itemize}

\subsection{Classic Physics Treatises \& Foundations}
The style cleanly formats landmark physics monographs and papers:
\begin{itemize}
  \item \textbf{Quantum Mechanics:} Dirac (\citeyear{dirac1958}, \parencite{dirac1958}), Wheeler and Zurek (\citeyear{wheeler1983}, \parencite{wheeler1983}), Bohm (\citeyear{bohm1983}, \parencite{bohm1983}), Planck (\citeyear{planck1922}, \parencite{planck1922}).
  \item \textbf{Relativity \& Classical Foundations:} Einstein (\citeyear{einstein1920}, \parencite{einstein1920}), Landau and Lifshitz (\citeyear{landau1976}, \parencite{landau1976}), Newton (\citeyear{newton1999}, \parencite{newton1999}), Lorentz (\citeyear{lorentz1952}, \parencite{lorentz1952}).
  \item \textbf{Astrophysics \& Gravitation:} Lynden-Bell (\citeyear{1967MNRAS.136..101L}, \parencite{1967MNRAS.136..101L}), Shakura and Sunyaev (\citeyear{shakura1973}, \parencite{shakura1973}), Sunyaev and Zeldovich (\citeyear{1970CoASP...2...66S}, \parencite{1970CoASP...2...66S}), Sample (\citeyear{sample2016}, \parencite{sample2016}).
  \item \textbf{Applied \& Materials Physics:} Sallo and Olsen (\citeyear{sallo1961}, \parencite{sallo1961}), Royall et al. (\citeyear{royall2013}, \parencite{royall2013}).
\end{itemize}

\section{Minimal Working Example (MWE)}

Here is a minimal template for your own thesis, paper, or report:

\begin{verbatim}
\documentclass[a4paper,11pt]{article}
\usepackage[british]{babel}
\usepackage{csquotes}
\usepackage[style=leeds-harvard]{biblatex}
\addbibresource{references.bib}
\usepackage[hidelinks]{hyperref}

\begin{document}

According to \textcite{shannon1948}, entropy measures uncertainty.
This theorem was verified across networks \parencite{tanenbaum2011}.

\printbibliography % Automatically produces section titled "References"

\end{document}
\end{verbatim}

\vspace{0.5em}
\noindent \textbf{The above MWE produces the following in-text output:}

According to \textcite{shannon1948}, entropy measures uncertainty.
This theorem was verified across networks \parencite{tanenbaum2011}.

\vspace{1em}
\noindent \textbf{And generates the following reference list entries:}

\begin{quote}
\textbf{References}

Shannon, C.E. 1948. A mathematical theory of communication. \emph{Bell System Technical Journal}. [Online]. \textbf{27}(3), pp.379-423. [Accessed 4 September 2024]. Available from: \url{https://doi.org/10.1002/j.1538-7305.1948.tb01338.x}

\vspace{0.6em}
Tanenbaum, A.S. and Wetherall, D.J. 2011. \emph{Computer networks}. 5th ed. Boston: Pearson Prentice Hall.
\end{quote}

\vspace{1.5em}
\noindent Below is the complete bibliography of all works cited across this test document:

\printbibliography

\end{document}
